\chapter{Conclusion}
\label{chap:conclusion}

\section{Research Questions Answered}

This chapter returns directly to the research questions posed at the start of the project.

\textbf{RQ1: How does backend choice affect checkpoint loading time for real LLM checkpoints?}

Backend choice affects loading time substantially, but not uniformly. At the Rust layer, \texttt{SafeTensors} and \texttt{ServerlessLLM} produce different backend rankings, and those rankings also change with checkpoint size. Smaller \texttt{SafeTensors} loads favour \texttt{sync}, while larger \texttt{SafeTensors} loads favour the final multi-worker \texttt{io\_uring} design. For \texttt{ServerlessLLM}, the strongest explicit backend is usually \texttt{async} or \texttt{io\_uring}, with \texttt{sync} rarely best.

\textbf{RQ2: How does checkpoint layout affect the relative performance of backends?}

Checkpoint layout changes the effective workload. Whole-file shard loading and range-heavy partitioned loading stress different parts of the implementation. \texttt{SafeTensors} rewards strong bulk shard parallelism; \texttt{ServerlessLLM} rewards good request grouping, coalescing, and range scheduling. The evaluation therefore shows that layout and loader must be considered together rather than in isolation.

\textbf{RQ3: How much overhead is introduced by the Python integration layer?}

The Python layer introduces significant overhead through runtime orchestration, tensor conversion, and object materialisation. That overhead is large enough to distort backend-level gains if left untreated. However, binding-specific optimisation reduced the public \texttt{default} path substantially on Qwen3-8B, showing that this overhead can be meaningfully reduced through engineering work at the binding layer.

\textbf{RQ4: To what extent do backend gains survive integration into vLLM?}

Backend gains survive into vLLM most clearly for model initialisation and TTFT. They matter much less for steady-state decode, where GPU compute dominates. The completed vLLM matrix confirms that loader choice still matters in a serving stack, but it also confirms that backend microbenchmarks alone are not enough to predict end-to-end behaviour.

\textbf{RQ5: What practical engineering lessons emerged?}

The most important lesson is that the architecture of a backend matters as much as the backend family. The early single-ring \texttt{io\_uring} prototype underperformed badly, but the final multi-worker design became competitive or best in several important regimes. The correct design target is therefore not one universal backend, but an adaptive, format-aware \texttt{default} policy backed by explicit override paths.

\section{Main Contributions}

The project makes four main contributions.

\begin{enumerate}
\item It builds a common Rust loading framework that allows direct comparison of \texttt{sync}, \texttt{async}, and \texttt{io\_uring} under shared APIs and shared formats.
\item It demonstrates that the performance ranking of these backends depends on workload structure rather than on a single universal rule.
\item It shows that binding-layer and serving-layer overheads can materially alter how backend gains appear in practice.
\item It produces an adaptive default policy that is substantially more defensible than either a fixed backend choice or the earlier synthetic cost-model approach.
\end{enumerate}

\section{Limitations}

The conclusions should still be interpreted with some care.

The experiments were run on one H100-based Linux environment, so the absolute timings are hardware-specific. The selector is much better than the earlier heuristic, but not perfect in every crossover case. The vLLM results are also partly conditioned on serving-engine settings such as memory budget and maximum model length. These limitations affect the exact quantitative boundaries between regimes, but they do not overturn the main qualitative result.

\section{Final Conclusion}

The final conclusion of the project is simple but stronger than the original hypothesis. The most useful way to think about LLM checkpoint loading is as an adaptive heuristic problem. The backend should not be fixed in advance; it should be selected from workload structure. In the clearest high-level form, the final evidence supports sync for smaller eager checkpoints and io\_uring for larger ones, with format-specific refinements around that rule.

That is the central outcome of this work. The strongest practical policy is an adaptive default, informed by real workload structure and validated across Rust, Python, and vLLM layers.
